\chapter*{УДИРТГАЛ}\addcontentsline{toc}{chapter}{УДИРТГАЛ}\markboth{УДИРТГАЛ}{}
\textbf{Сэдвийн үндэслэл.} Улаанбаатар хотод бүртгэлтэй тээврийн хэрэгслийн тоо 2025 оны байдлаар 730 мянгаар хязгаарлагдах шийдвэр гарсан бөгөөд зогсоолын хүрэлцээ, төлбөрийн автоматжуулалт нь хотын дэд бүтцийн тулгамдсан асуудлын нэг болоод байна~\cite{asianews2025}. Зогсоолын хаалгыг автоматаар нээх, төлбөрийг машины дугаараар тооцох, орох-гарах үйл явдлыг тасалбаргүйгээр тулгах зэрэг үйлдлийн суурь нь улсын дугаарыг камерын кадраас автоматаар таних (Automatic License Plate Recognition, ALPR) технологи юм. Олон улсын арилжааны ALPR шийдлүүд нь латин цагаан толгой эсвэл хятад дугаарт тохируулагдсан бол Монгол улсын дугаар нь орос цагаан толгойд байдаггүй Ө, Ү үсэг бүхий монгол кирилл үсгээр бичигддэг, нэг мөрт ба хоёр мөрт гэсэн хоёр физик хэлбэртэй~\cite{wikimn}. Иймээс бэлэн OCR загварууд монгол дугаарыг шууд уншиж чаддаггүй, дугаарын дүрмийг мэддэггүй тул латин--кирилл ижил хэлбэрт үсгийг андуурдаг.

\textbf{Сэдвийг сонгосон шалтгаан.} Нэгдүгээрт, монгол дугаарт зориулсан нээлттэй, текстээр шошголсон өгөгдлийн сан, хэвлэгдсэн судалгааны ажил байхгүй тул өгөгдлийг нэгтгэх, үнэлгээний протокол тогтоох нь өөрөө шинэлэг. Хоёрдугаарт, YOLO~\cite{yolo11} хэмээх хөнгөн илрүүлэгч, CRNN+CTC~\cite{shi2017crnn,graves2006ctc} дарааллын уншигч зэрэг сүүлийн үеийн загварууд нь нэг кадрыг хэдхэн миллисекундэд боловсруулж чадах болсноор зогсоолын хаалга дээр бодит цагт ажиллах боломж бүрдсэн. Гуравдугаарт, нэг кадрын алдааг олон кадрын саналаар засах арга нь видео ALPR-д 10--15 нэгж хувийн ахиц өгдөг нь судалгаагаар батлагдсан~\cite{laroca2018,wang2026mt} тул үүнийг монгол дугаарт хэрэгжүүлэх нь практик ач холбогдолтой.

\textbf{Судалгааны зорилго.} Зогсоолын системийн мэдрэхүйн модуль болох, Монгол улсын дугаарыг камерын кадраас 95 хувиас дээш бүрэн зөв уншдаг, нэг кадрыг 100 миллисекундээс бага хугацаанд боловсруулдаг, олон кадраар баталгаажуулдаг гүн сургалтын загвар бүтээхэд оршино.

\textbf{Судалгааны зорилтууд.} Зорилгод хүрэхийн тулд дараах зорилтуудыг дэвшүүлэв:
\begin{enumerate}[leftmargin=2.5em]
\item Дугаар таних олон улсын судалгаа, монгол дугаарын бүтэц, нээлттэй өгөгдлийн эх үүсвэрүүдийг судалж, цөөн өгөгдөлтэй нөхцөлд тохирох хандлагыг сонгох;
\item Олдсон бүх эх үүсвэрийг (GitHub, Roboflow, HuggingFace, MLUB тэмцээн) нэгтгэж, тэмдэгтийн хайрцгаас дугаарын текст сэргээх, дугаарын давхцалгүй хуваалт хийх, синтетик дугаар үүсгэх боловсруулалтын урсгал бүтээх;
\item YOLO11n дугаар илрүүлэгчийг хоёр үе шаттай сургаж mAP$_{50}\ge0.95$-д хүргэх;
\item Монгол кирилл үсгийн бүрэн олонлог дээр CRNN+CTC уншигчийг синтетик ба бодит өгөгдлөөр үе шаттай сургаж, нэг ба хоёр мөрт дугаарыг хоёуланг нь уншдаг болгох;
\item Дугаарын дүрэм, аймгийн кодын хүснэгт болон итгэлээр жинлэсэн тэмдэгт тус бүрийн саналд суурилсан олон кадрын баталгаажуулалт хэрэгжүүлэх;
\item Үр дүнг хэмжүүрээр (mAP, бүтэн дугаарын нарийвчлал, CER, хугацаа) үнэлж, алдааны шинжилгээ хийж, туршилтын интерфэйс бүтээн, дараагийн шатны ажлыг тодорхойлох.
\end{enumerate}

\textbf{Судалгааны хүрээ.} Энэ тайланд тусгасан ажил нь зогсож буй эсвэл удаан хөдөлж буй тээврийн хэрэгслийн дугаарыг өдрийн гэрэлд таних асуудалд хамаарна. Зогсоолын бүрэн систем (орох-гарах бүртгэлийн сан, төлбөр, хаалганы удирдлага) болон шөнийн, цас борооны нөхцөл дараагийн шатанд хамаарна. Дипломат (улаан, ДК), цэргийн, чиргүүлийн тусгай дугаар энэ ажилд ороогүй.

\textbf{Явцын байдал.} Ажлыг 2026 оны 10 дугаар сарын 14-нд эхлүүлж, нээлттэй эх үүсвэрээс 68 бүтэн дүрстэй зураг, 7392 бодит дугаарын тайралт, 8823 олон улсын зураг, 3 бичлэг цуглуулж, 90\,021 синтетик мөр зураг үүсгэв. Илрүүлэгч монгол шалгалтын олонлог дээр mAP$_{50}=0.992$, уншигч 1491 бодит дугаар дээр 96.8 хувийн бүрэн зөв уншилт (CER 0.65 хувь), бүхэл урсгал нэг 1080p кадрыг 7.3~мс-д боловсруулж байна. Бүтэн дүрстэй 19 дугаарын жижиг тест дээр төгсгөл хүртэлх нарийвчлал 73.7 хувь (9 кадрын саналаар 78.9 хувь) байгаа нь домэйн зөрүү ба ховор үсгийн асуудлыг илтгэж байгааг 5 дугаар бүлэгт хэлэлцэв.

\textbf{Тайлангийн бүтэц.} 1 дүгээр бүлэгт дугаар таних судалгааны тойм, 2 дугаар бүлэгт өгөгдлийн эх үүсвэр, боловсруулалт, синтетик үүсгэлт, 3 дугаар бүлэгт таних системийн арга зүй ба сургалтын тохиргоо, 4 дүгээр бүлэгт туршилт, үр дүн, алдааны шинжилгээ, 5 дугаар бүлэгт хэлэлцүүлэг, хязгаарлалт, цаашдын ажлыг тус тус тусгав.
